
\section{Results}
\label{sec:results}

\subsection{AI Forecasting Performance}

Table~\ref{tab:modelarena} presents comparative prequential accuracy across all nine
estimators plus a naive Persistence baseline on the 500-sample controlled synthetic
benchmark. The Persistence baseline ($\hat{y}_{t+h}=y_t$) establishes the minimum
performance floor; all incremental estimators substantially outperform it. The blended
ensemble \textbf{M8 (SGD+PA, 0.6/0.4)} achieved the best performance across all metrics,
with MAE of \textbf{10.616~\ugsm}, RMSE of \textbf{17.431~\ugsm}, and $R^2=\textbf{0.7334}$.
The first sample (k=1) is excluded from prequential scoring (no prior prediction),
yielding $N=499$ scored samples.

\begin{table}[htbp]
\centering
\caption{Prequential accuracy and hazard metrics for the Persistence baseline,
  all nine Model Arena estimators (M1--M9) on the controlled synthetic benchmark
  ($N=499$ scored samples). \textbf{Bold}: champion M8. $^\dagger$Persistence values
  are stream-statistic estimates (not prequentially evaluated).}
\label{tab:modelarena}
\footnotesize
\setlength{\tabcolsep}{4pt}
\begin{tabular}{@{}clccccc@{}}
\toprule
\textbf{ID} & \textbf{Model} & \textbf{MAE (\ugsm)} & \textbf{RMSE (\ugsm)} &
\textbf{$R^2$} & \textbf{Recall} & \textbf{$F_1$} \\
\midrule
-- & Persistence$^\dagger$ ($\hat{y}=y_{t-1}$) & 18.9 & 27.1 & 0.42 & 0.804 & 0.738 \\
\midrule
M1 & SGD -- $L_2$ Ridge        & 11.200 & 18.400 & 0.712 & 0.951 & 0.925 \\
M2 & SGD -- $L_1$ Lasso        & 12.450 & 19.820 & 0.689 & 0.948 & 0.919 \\
M3 & SGD -- ElasticNet         & 11.900 & 19.150 & 0.701 & 0.955 & 0.926 \\
M4 & SGD -- Huber ($\delta$=1.35) & 10.950 & 17.920 & 0.724 & 0.961 & 0.935 \\
M5 & SGD -- SVR ($\varepsilon$-insens.) & 12.600 & 20.050 & 0.684 & 0.942 & 0.913 \\
M6 & Passive-Aggressive $C$=1.0 & 12.800 & 20.100 & 0.681 & 0.970 & 0.927 \\
M7 & Passive-Aggressive $C$=0.1 & 13.400 & 21.200 & 0.665 & 0.958 & 0.914 \\
\textbf{M8} & \textbf{SGD+PA Blend 0.6/0.4~[Champion]} & \textbf{10.616} & \textbf{17.431} & \textbf{0.733} & \textbf{0.965} & \textbf{0.940} \\
M9 & EMA Baseline (non-ML)     & 15.800 & 24.300 & 0.592 & 0.890 & 0.854 \\
\bottomrule
\end{tabular}
\end{table}

\begin{figure}[htbp]
\centering
\includegraphics[width=0.92\linewidth]{model_arena_comparison.png}
\caption{Comparative prequential MAE and RMSE across all nine Model Arena estimators (M1--M9)
  on the controlled synthetic benchmark ($N=499$ scored samples; test-then-train protocol).
  Champion M8 (SGD+PA Blend) achieves the lowest error (MAE=10.616~\ugsm, RMSE=17.431~\ugsm);
  M9 provides the non-ML EMA reference floor; the dashed horizontal line indicates the naive
  Persistence baseline ($\hat{y}_{t+h}=y_t$, MAE=18.90~\ugsm).}
\label{fig:arena-results}
\end{figure}

On the Bangladesh public-dataset-derived stream ($N=500$), the Model Arena achieved
aggregate $R^2=0.8853$. M6 (PA, $C=1.0$) achieved the lowest standalone MAE
(8.552~\ugsm), while the blended M8 remained competitive (MAE 9.100~\ugsm), suggesting
reasonable transfer to ambient atmospheric dynamics under the feature mapping used.

\begin{figure}[htbp]
\centering
\includegraphics[width=0.65\linewidth]{dashboard_real_top.jpg}\\[3pt]
\includegraphics[width=0.65\linewidth]{dashboard_real_bottom.jpg}
\caption{NirmalNode live operator web dashboard during active monitoring.
  \textit{Top panel:} AQI ring gauge (AQI~80, PM$_{2.5}=25.8$~\ugsm), 1-hour and
  2-hour indicative forecasts, XAI/LFC explanation panel, and real-time Model Arena
  prequential leaderboard (champion M8-BLEND highlighted, MAE~5.300 at sample~79).
  \textit{Bottom panel:} Particulate matter channels (PM$_{1.0}$, PM$_{2.5}$,
  PM$_{10}$, optical bin counts), multi-gas sensor ADC readings (MQ-135, MQ-136,
  MiCS RED/NOX, MP-135), live PM$_{2.5}$ vs.\ AI forecast chart, and prequential
  online MAE/RMSE convergence plot.}
\label{fig:dashboard}
\end{figure}

\subsection{Hazard Detection}

Table~\ref{tab:confusion} presents binary confusion matrices for both evaluation streams.
On the controlled benchmark, M8 achieves hazard Precision~0.916, Recall~0.965,
and $F_1=0.940$, with 4 missed hazard events (FN=4, 0.8\%) out of 500. On the
Bangladesh public-dataset-derived stream, the model achieves zero false negatives
(Recall~1.000, $F_1=0.992$) in $N=500$ evaluated samples; the exact 95\% binomial
confidence interval for Recall on this stream is [0.990, 1.000], reflecting the
limitation of a single evaluation run. AUC-ROC~0.985 and PR-AUC~0.962 on this stream.

\begin{table}[htbp]
\centering
\caption{Confusion matrices and classification performance for the predictive purifier
  trigger ($\tau=35.5$~\ugsm). FN~=~missed hazard event.}
\label{tab:confusion}
\footnotesize
\begin{tabular}{@{}lcccccccc@{}}
\toprule
\textbf{Stream} & \textbf{TP} & \textbf{FP} & \textbf{TN} & \textbf{FN} &
\textbf{Precision} & \textbf{Recall} & \textbf{$F_1$} & \textbf{Accuracy} \\
\midrule
Controlled synth. ($N=500$)          & 109 & 10 & 377 & 4  & 0.916 & 0.965 & 0.940 & 97.2\% \\
Bangladesh pub.-data ($N=500$) & 381 & 6  & 113 & \textbf{0} & 0.984 & \textbf{1.000} & \textbf{0.992} & 98.8\% \\
\bottomrule
\end{tabular}
\end{table}

\begin{figure}[htbp]
\centering
\begin{subfigure}[b]{0.48\linewidth}
  \includegraphics[width=\linewidth]{confusion_matrix.png}
  \caption{Confusion matrix (controlled stream).}
\end{subfigure}
\hfill
\begin{subfigure}[b]{0.48\linewidth}
  \includegraphics[width=\linewidth]{roc_pr_curves.png}
  \caption{ROC (AUC=0.985) and PR (AUC=0.962).}
\end{subfigure}
\caption{Binary classification diagnostics for the predictive purifier trigger.}
\label{fig:roc}
\end{figure}

\subsection{Online Adaptation and Error Convergence}

Figure~\ref{fig:convergence} shows the prequential rolling MAE across the full 500-sample
lifespan in three operationally meaningful phases:

\begin{figure}[htbp]
\centering
\begin{subfigure}[b]{0.82\linewidth}
  \centering
  \includegraphics[width=\linewidth]{prequential_error_plot.jpg}
  \caption{Rolling prequential MAE and RMSE across Cold-Start (1--50), Adaptation (50--200), and Convergence (200--500) phases.}
  \label{fig:prequential-error}
\end{subfigure}\\[6pt]
\begin{subfigure}[b]{0.58\linewidth}
  \centering
  \includegraphics[width=\linewidth]{residual_distribution.png}
  \caption{Residual distribution and Q--Q normality plot ($\mu=-0.12$~\ugsm, $\sigma=17.4$~\ugsm).}
  \label{fig:residual-qq}
\end{subfigure}
\caption{Online adaptation evidence and residual error diagnostics for champion model M8.}
\label{fig:convergence}
\end{figure}

\textbf{Cold-Start (samples 1--50):} MAE drops from \textbf{33.4~\ugsm} to 13.9~\ugsm{}
as the ensemble absorbs site-specific data.
\textbf{Adaptation (50--200):} MAE declines to 11.4~\ugsm{} as feature covariance stabilizes.
\textbf{Convergence (200--500):} MAE stabilizes at \textbf{10.616~\ugsm}
(RMSE~17.431~\ugsm), a \textbf{68.2\% reduction} from cold-start MAE via streaming
adaptation without manual retraining. The residual mean is close to zero
($\mu=-0.12$~\ugsm), indicating limited average bias; approximate residual normality is
assessed qualitatively via the Q-Q plot (Fig.~\ref{fig:convergence}b).

\subsection{Explainability via Linear Feature Contribution (XAI)}

Because M8 blends regularized linear estimators, each prediction decomposes exactly into
per-feature Linear Feature Contribution (LFC) terms under standardized inputs:
\begin{equation}
\phi_i = w_{\mathrm{blend},i}\cdot\frac{x_i-\mu_i}{\sigma_i},\qquad
\hat{y} = b_{\mathrm{blend}}+\sum_{i=1}^{14}\phi_i
\label{eq:lfc}
\end{equation}

\noindent where $w_{\mathrm{blend},i}$ is the blended model coefficient, $\mu_i$ and
$\sigma_i$ are the running online scaler statistics. This provides fully interpretable,
instance-level attribution without requiring a separate explanation model.

\begin{figure}[htbp]
\centering
\includegraphics[width=0.60\linewidth]{xai_feature_importance.png}
\caption{Global feature importance ranking of champion model M8.
  PM$_{2.5}$-related features (EMA, current, lag-1 delta) dominate
  $>$64\% of total decision weight, supporting physically interpretable learning.}
\label{fig:xai}
\end{figure}

For a high-pollution instance (PM$_{2.5}=78$~\ugsm), the lag-1 delta applied restorative
damping ($\phi=-27.80$~\ugsm), while the PM$_{2.5}$ EMA ($\phi=+15.47$), optical counts
($+14.65$), PM$_{10}$ ($+14.10$), and current PM$_{2.5}$ ($+13.76$) elevated the forecast
toward the hazard range. This feature alignment with atmospheric aerosol transport physics
supports operational interpretability and physical plausibility assessment of model behaviour.

\subsection{Physical Purification Efficiency}

Table~\ref{tab:purif-pm} reports quantitative single-pass PM removal measured by
paired PMS5003 sensors in the aerosol bench trial. Gas-phase responses are reported
separately in Table~\ref{tab:purif-gas} as qualitative indicators only.

\begin{table}[htbp]
\centering
\caption{Quantitative PM single-pass removal (bench trial, incense smoke, 0.85~m$^3$
  chamber, single measurement run). PMS5003 is not a reference-grade instrument;
  results should be interpreted as bench prototype estimates.}
\label{tab:purif-pm}
\footnotesize
\begin{tabular}{@{}llccc@{}}
\toprule
\textbf{Pollutant} & \textbf{Stage} & \textbf{$C_\mathrm{in}$ (\ugsm)} &
\textbf{$C_\mathrm{out}$ (\ugsm)} & \textbf{$\eta_p$ (\%)} \\
\midrule
PM$_{1.0}$ & HEPA-H13         & 42  & 1.8 & \textbf{95.7} \\
PM$_{2.5}$ & HEPA-H13         & 67  & 2.9 & \textbf{95.7} \\
PM$_{10}$  & Cyclone+HEPA     & 98  & 1.4 & \textbf{98.6} \\
\bottomrule
\end{tabular}
\end{table}

\begin{table}[htbp]
\centering
\caption{\textbf{Qualitative gas-phase observations only} --- MOS differential
  ADC readings before/after activated-carbon stage. These values are \emph{not}
  calibrated gas concentrations and are not comparable to quantitative PM results.
  Reference-grade electrochemical analysis was not available.}
\label{tab:purif-gas}
\footnotesize
\begin{tabular}{@{}llcc@{}}
\toprule
\textbf{Gas} & \textbf{Stage} & \textbf{Observation} & \textbf{Approx. ADC change} \\
\midrule
VOCs   & Activated carbon & Reduced ADC response & $\sim$60--70\% decrease (qualitative) \\
H$_2$S & Activated carbon & Below sensor noise floor & $>$50\% decrease (qualitative) \\
\bottomrule
\end{tabular}
\end{table}

In the aerosol bench trial, the predictive engine energized the blower \textbf{4--6 seconds
prior} to the monitored concentration crossing $\tau=35.5$~\ugsm, establishing pre-actuation
airflow before peak plume arrival. This single-trial observation should be interpreted as
a bench-scale feasibility indication; repeated trials with statistical characterization
are needed to fully establish the lead-time distribution. Measured PM$_{2.5}$ bench
removal (95.7\%) is consistent with independent HEPA literature: Lu~\etal{}~\citep{lu2023realworld}
reported 75--91\% and Cox~\etal{}~\citep{cox2018effectiveness} reported 73--95\% in
controlled settings.

\subsection{Power and Prototype Cost}

\begin{table}[htbp]
\centering
\caption{Measured power profile and prototype Bill of Materials (BOM) at retail
  Dhaka pricing (BDT~110/USD, August~2025). Prototype component cost only; does not
  include manufacturing, certification, labour, or logistics.}
\label{tab:power}
\footnotesize
\setlength{\tabcolsep}{4pt}
\begin{tabular}{@{}p{5.2cm}p{3.8cm}cc@{}}
\toprule
\multicolumn{4}{@{}l}{\textit{(a) Power profile}} \\
\midrule
\textbf{Subsystem} & \textbf{State} & \textbf{Power (W)} \\
\cmidrule{1-3}
ESP32 + sensors + Wi-Fi & TX burst   & 2.6 \\
12~V DC blower          & Fan ON     & 6.0 \\
\textbf{System (standby)}\quad & Fan OFF   & \textbf{1.8} \\
\textbf{System (peak)}  & Fan ON + Wi-Fi & \textbf{8.6} \\
\midrule
\multicolumn{4}{@{}l}{\textit{(b) Prototype BOM (selected items)}} \\
\midrule
\textbf{Component} & \textbf{Spec.} & \textbf{BDT} & \textbf{USD} \\
\cmidrule{1-4}
ESP32-WROOM-32      & Dev module     & 350  & \$3.18 \\
PMS5003 OPC         & Plantower      & 2200 & \$20.00 \\
MiCS-4514           & SGX dual gas   & 1650 & \$15.00 \\
MQ135/MQ136/MP135   & 3$\times$ MOS  & 580  & \$5.27 \\
BME280              & Bosch T/RH/P   & 280  & \$2.55 \\
Relay + blower      & 12~V DC        & 490  & \$4.45 \\
Cyclone+HEPA+carbon & 3-stage train  & 1250 & \$11.36 \\
Enclosure/wiring    & Acrylic/PVC    & 600  & \$5.45 \\
\textbf{Total}      &                & \textbf{7,400} & \textbf{\$67.27} \\
\bottomrule
\end{tabular}
\end{table}

Under a 22.6\% duty cycle (proportion of samples above $\tau$), estimated battery life
from a 55.5~Wh pack is approximately 17~hours, assuming 85\% converter efficiency and
full usable capacity (calculation: $55.5 \times 0.85 / (0.226 \times 8.6 + 0.774 \times 1.8)
\approx 17.1$~h), sufficient for a typical industrial double-shift.
